\prefacesection{Declaration of Academic Achievement}

The student will declare his/her research contribution and, as appropriate, 
those of colleagues or other contributors to the contents of the thesis. 
\\~\\
{\color{red}Dr.~Zhao developed most of the mathematics for the formulation of the manifold $\M_C$, including the subspace $V$ \ref{eq:subspace:v}, and Theorem \ref{thm:Theorem:1.3} and its proof. Dr.~Zhao also developed the mathematics for much of the Euler-Voigt optimization problem described in Chapter \ref{chp:optimization:problem}, including the objective functional \ref{eq:objective:functional}, the linearization of the Euler-Voigt equations \ref{eq:euler:voigt:linearization}, the adjoint system \ref{eq:euler:voigt:adjoint}, and the gradient of the objective functional \ref{eq:gradient}. The Riemannian conjugate gradient method for the Euler-Voigt optimization problem in Section \ref{sec:Riemannian Conjugate Gradient Method} was developed by Dr.~Zhao and Dr.~Protas. This includes the tangent space \ref{eq:tangent:space}, The projection operator \ref{eq:projection:operator}, the search direction \ref{eq:search:direction}, and step size \ref{eq:step:size}. The retraction operator \ref{eq:retraction:operator}, and the vector transport operation \ref{eq:vector:transport} were originally developed by Dr.~Zhao and Dr.~Protas but I modified to suit this project.
\\~\\
The size of the manifold $\M_C$ constraint $C$, and the derivations in Appendix E and F are my work. The results in Appendix A, B, C, and D are my re-derivations of Dr.~Zhao's original mathematical results. The derivation in Appendix G and H are my re-derivations of the well known Euler and Euler-Voigt scaling properties. The determination of the numerical model constants, time step size and $\sigma$, were also my work. The results and conclusions in Chapters 5 and 6 are my work. For the numerical models, Dr.~Zhao created the original code base which I modified to suit this project. The Python code used to perform data analysis and produce the graphs is entirely my own work. }